% =====================================================================
% §1.7 — Preparación de sustrato e inoculación
% =====================================================================
\section{Preparación de sustrato e inoculación}
\label{sec:preparacion}

Subsistema \emph{aguas arriba} del rack: prepara y esteriliza el
sustrato, propaga el micelio en grano (\emph{grain spawn}) y arma las
bolsas inoculadas que se incubarán y trasladarán a los contenedores de
fructificación (\S\ref{sec:contenedores}).

\textbf{Nota.} Esta sección reemplaza el supuesto inicial del brief de
``recibir bolsas ya inoculadas''. El operador realiza el proceso
completo desde sustrato bruto.

% ---------------------------------------------------------------------
\subsection{Esterilización por olla de presión}
\label{sec:preparacion:autoclave}

Se usa una olla de presión doméstica de \paramVolOllaPresion\,L como
autoclave de bajo costo. Parámetros de operación:

\begin{table}[H]
  \centering
  \caption{Ciclos de esterilización en olla de presión.}
  \label{tab:autoclave}
  \begin{tabularx}{\linewidth}{Xcc}
    \toprule
    \textbf{Material} & \textbf{Presión [psi]} & \textbf{Tiempo [min]} \\
    \midrule
    Grano para spawn (frasco 1\,L)               & \paramPresionEster & \paramTiempoEsterGrano \\
    Sustrato en bolsa (paja, pellets, mezclas)   & \paramPresionEster & \paramTiempoEsterSustrato \\
    Material de inoculación (bisturí, jeringas)  & \paramPresionEster & 30 \\
    \bottomrule
  \end{tabularx}
\end{table}

\paragraph{Validación de ciclo.} El tiempo \emph{efectivo} comienza
cuando la válvula libera vapor sostenido; el termómetro de aguja
montado en la tapa permite verificación cruzada. Para sustratos densos
o cargas mayores se incrementa el tiempo en 30\,min por cada nivel
adicional de carga.

\paragraph{Enfriado.} La bolsa o frasco esterilizado debe enfriarse a
$T \leq 30\,\si{\celsius}$ antes de inoculación (típicamente 8--12\,h en
ambiente cerrado). \textbf{No abrir} la bolsa antes de inoculación.

\subsection{Still Air Box (SAB)}
\label{sec:preparacion:sab}

Cámara de aire estático construida con tote hermético de 60--80\,L
adicional, con dos perforaciones circulares (diámetro $\approx$
$\sfrac{1}{2}$ antebrazo) en la pared corta y manga elástica opcional.
Procedimiento de uso:

\begin{enumerate}[leftmargin=*,itemsep=0pt]
  \item Limpieza con peróxido 3\,\% + alcohol isopropílico 70\,\%
        interior y exterior.
  \item Reposo cerrado 5--10\,min para sedimentación de aerosoles.
  \item Introducir bolsa esterilizada, jeringa de spawn, bisturí
        flameado, alcohol en spray.
  \item Inocular minimizando movimientos amplios; sellar la perforación
        del filtro con cinta micropore inmediatamente tras inyectar.
\end{enumerate}

\pendientecurso{Validar dimensiones exactas del SAB y procedimiento
contra recomendaciones del curso. Considerar upgrade a flow hood HEPA
casero si la tasa de contaminación supera 10\,\% en los primeros 5
lotes.}

\subsection{Producción de grain spawn}
\label{sec:preparacion:spawn}

El grain spawn (micelio sobre grano) es el inóculo intermedio entre el
cultivo madre (jeringa de esporas o cultivo en placa) y el sustrato
final. Proceso:

\begin{enumerate}[leftmargin=*,itemsep=0pt]
  \item Enjuagar y remojar grano (quinoa, trigo o mijo) 12--24\,h.
  \item Escurrir y secar superficialmente.
  \item Frasco 1\,L cargado a $\sfrac{2}{3}$; tapa con perforación
        para intercambio gaseoso cubierta con cinta micropore.
  \item Esterilizar en olla \paramVolOllaPresion\,L: \paramTiempoEsterGrano\,min
        a \paramPresionEster\,psi.
  \item Enfriar a $T \leq 30\,\si{\celsius}$ (mínimo 8\,h).
  \item Inocular con jeringa de esporas o transferir trozos de placa
        en SAB.
  \item Incubar en oscuridad a 22--24\,\si{\celsius} hasta colonización
        completa (10--21\,días según especie).
  \item Agitar el frasco una vez a la mitad de la colonización para
        distribuir micelio.
\end{enumerate}

\subsection{Armado de bolsas de sustrato inoculado}
\label{sec:preparacion:bolsas}

\begin{enumerate}[leftmargin=*,itemsep=0pt]
  \item Hidratar sustrato a campo (HR $\approx 65$\,\%) y agregar
        suplemento (yeso, salvado) según especie.
  \item Embolsar 1{,}5--2{,}5\,kg por bolsa autoclavable con parche
        filtro 0{,}2\,\si{\micro\meter}.
  \item Esterilizar a \paramPresionEster\,psi durante
        \paramTiempoEsterSustrato\,min.
  \item Enfriar.
  \item Inocular en SAB con $\approx 5$--$10$\,\% del peso húmedo en
        grain spawn.
  \item Cerrar bolsa con sello impulso o liga doblada.
  \item Etiquetar (lote, especie, fecha) y trasladar a incubación.
\end{enumerate}

\subsection{Control de contaminación}
\label{sec:preparacion:contam}

Registro obligatorio por lote: número de bolsas inoculadas, número de
bolsas con contaminación detectada (Trichoderma verde, mohos negros,
bacterial sour), fecha de detección. Métricas objetivo:

\begin{itemize}[leftmargin=*,itemsep=0pt]
  \item Tasa de contaminación $<$10\,\% en los primeros 5 lotes.
  \item Tasa $<$5\,\% en régimen establecido.
  \item Una tasa $>$25\,\% indica revisar: tiempo de esterilización,
        sello del filtro de la bolsa, técnica en SAB o calidad del
        cultivo madre.
\end{itemize}

\subsection{Hooks para la base de datos (Tópico~7)}
\label{sec:preparacion:hooks}

Campos adicionales en la tabla \texttt{lotes}:

\begin{itemize}[leftmargin=*,itemsep=0pt]
  \item \texttt{sustrato\_tipo}, \texttt{sustrato\_masa\_hidratada\_kg}
  \item \texttt{esteriliz\_temp\_C}, \texttt{esteriliz\_tiempo\_min},
        \texttt{esteriliz\_presion\_psi}
  \item \texttt{spawn\_origen} (jeringa lote / placa propia)
  \item \texttt{spawn\_porcentaje} (\% sobre masa húmeda)
  \item \texttt{contaminacion\_detectada\_bool},
        \texttt{contaminacion\_tipo} (enum)
  \item \texttt{fecha\_inoculacion}, \texttt{fecha\_colonizacion\_completa}
\end{itemize}

\subsection{Criterios de aceptación}
\label{sec:preparacion:aceptacion}

\begin{itemize}[leftmargin=*,itemsep=0pt]
  \item Olla de presión instalada, ciclos verificados con termómetro
        independiente.
  \item SAB armado y sanitizado.
  \item Primer lote de grain spawn colonizado a $>$90\,\%.
  \item Tasa de contaminación documentada para cada lote.
  \item Campos en \texttt{lotes} reflejados en el schema del Tópico~7.
\end{itemize}
